\section{Conclusions}
\label{sec:concs}

In this work we investigated the computational expressivity of fixed point logics.
Our main contribution is a characterisation of the functions representable in the systems $\muLJ$ and $\cmuLJ$ as the functions provably recursive in the second-order theory $\PSCA$.
This extends the tradition of such correspondences between arithmetic theories an type systems, in particular including between $\PA$ and system $\T$ due to G\"odel \cite{goedel1958bisher}, and between $\SOPA$ and $\F$ due to Girard \cite{girard1972systemF}.
In a sense ours is a somewhat intuitive result in light of M\"ollerfeld's work \cite{Moellerfeld02:phd-thesis}, identifying the latter's arithmetical theorems with the extension of $\ACA$ by least and greatest general fixed points.


Our characterisation also applies to aforementioned (circular) systems of linear logic, namely $\muMALL$ and its circular counterpart, from \cite{BaeldeMiller07,Baelde12,BaeldeDS16Infinitaryprooftheory,EhrJaf21:cat-models-muLL,EhrJafSau21:totality-CmuMALL,BDKS22:bouncing-threads,DeSau19:infinets1,DePelSau:infinets2,DeJafSau22:phase-semantics}.
This result was given the preliminary conference version of this paper \cite{CD23:muLJ-lics}, but will be expanded upon in a separate full paper.


Referring to Rathjen's ordinal notation system in \cite{Rathjen95:recent-advances}, this means that all these systems represent just the functions computable by recursion on ordinals of $\mathfrak T[\theta = \omega]$.
To this end we used a range of techniques from proof theory, reverse mathematics, higher-order computability and metamathematics.
This contribution settles the question of computational expressivity of (circular) systems with fixed points, cf.~Question~\ref{question:main}.

In future work it would be interesting to investigate the computational expressivity of systems with only \emph{strictly} positive fixed points, where $\mu,\nu$ may only bind variables that are \emph{never} under the left of $\arrow$.
We suspect that such systems are computationally weaker than those we have considered here.
On the arithmetical side, it would be interesting to investigate the \emph{higher-order} reverse mathematics of the Knaster-Tarski theorem, cf.~\cite{SatoYamazaki17:rev-math-fixed-point-thms-with-kt} and the present work.

\comment{
\begin{itemize}
    \item what about strict positivity? conjecture weaker.
    \item we could first translate into a well-founded system, but this will not (immediately) give us effective ordinal bounds for formalisation. (Actually this is an interesting idea)
\end{itemize}
}